\section*{Appendix}


\subsection*{Data Dictionary}

% Table 1: Cell Level
\begin{table}[H]
\centering
\caption{Cell Level Variables}
\label{tab:cell_level}
\begin{tabular}{lll}
\toprule
Variable & Description & Type \\
\midrule
\texttt{cycle\_life} & Number of cycles until EoL & Float \\
\texttt{charge\_policy} & Charging protocol & String \\
\texttt{summary} & Metadata for each cycle & Dictionary \\
\texttt{cycles} & Time series of measurements & Dictionary \\
\bottomrule
\end{tabular}
\end{table}

% Table 2: Summary Level
\begin{table}[H]
\centering
\caption{Summary Level (Cycle-by-Cycle Metadata)}
\label{tab:summary_level}
\begin{tabular}{lll}
\toprule
Variable & Description & Type \\
\midrule
\texttt{IR} & Internal resistance & Numeric array \\
\texttt{QC} & Charge capacity & Numeric array \\
\texttt{QD} & Discharge capacity & Numeric array \\
\texttt{Tavg} & Average temperature & Numeric array \\
\texttt{Tmin} & Minimum temperature & Numeric array \\
\texttt{Tmax} & Maximum temperature & Numeric array \\
\texttt{chargetime} & Charge time & Numeric array \\
\texttt{cycle} & Cycle number & Numeric array \\
\bottomrule
\end{tabular}
\end{table}

% Table 3: Cycle Level
\begin{table}[H]
\centering
\caption{Cycle Level (Intra-cycle Time Series)}
\label{tab:cycle_level}
\begin{tabular}{lll}
\toprule
Variable & Description & Type \\
\midrule
\texttt{I} & Current & Numeric array \\
\texttt{Qc} & Charge capacity & Numeric array \\
\texttt{Qd} & Discharge capacity & Numeric array \\
\texttt{Qdlin} & Linear interpolated discharge capacity & Numeric array \\
\texttt{T} & Temperature & Numeric array \\
\texttt{Tdlin} & Linear interpolated temperature & Numeric array \\
\texttt{V} & Voltage & Numeric array \\
\texttt{dQdV} & Derived vectors of discharge & Numeric array \\
\texttt{t} & Time & Numeric array \\
\bottomrule
\end{tabular}
\end{table}

\subsection*{Model Variables and Features}


\begin{table}[H]
\centering
\caption{Variables}
\begin{tabular}{ll}
\toprule
Symbol & Description \\
\midrule
$j = 1, \dots, J$ & Cycling condition group index \\
$i = 1, \dots, n_j$ & Cell index within group $j$ \\
$y_{ji}$ & EoL time (days) for cell $i$ in group $j$ \\
$\boldsymbol{\theta}_j \in \mathbb{R}^4$ & Regression coefficients for group $j$ \\
$\sigma^2_j$ & Variance of group $j$ model \\
$\boldsymbol{\gamma} \in \mathbb{R}^{4 \times 2}$ & Hyperparameter matrix linking cycling conditions to group coefficients \\
\bottomrule
\end{tabular}
\end{table}

\begin{table}[H]
\centering
\caption{Feature vectors}
\begin{tabular}{lll}
\toprule
Symbol & Components & Level \\
\midrule
$\mathbf{x}_{ji} \in \mathbb{R}^4$ & $(1, F1_{ji}, F2_{ji}, F3_{ji})^\top$ & Individual \\
$\mathbf{g}_j \in \mathbb{R}^2$ & $(1, \bar{g}_j)^\top$ & Group \\
\bottomrule
\end{tabular}
\end{table}

\begin{table}[H]
\centering
\caption{Individual-level features (computed from first 100 cycles)}
\begin{tabular}{ll}
\toprule
Symbol & Description \\
\midrule
$F1_{ji}$ & Variance of $\Delta Q(V)$ between cycles 10 and 100 \\
$F2_{ji}$ & Minimum of $\Delta Q(V)$ over the same interval \\
$F3_{ji}$ & Discharge capacity at cycle 2 \\
\bottomrule
\end{tabular}
\end{table}
\footnote{$\Delta Q(V)$ denotes the discharge curve difference as a function of voltage, following Severson et al. (2019).}
